\subsection{Our Results}

Our first result gives a lower bound on the degree needed to refute
$\circuit_s(f)$ in SoS.  This lower bound is very general and in fact
applies to \emph{every} Boolean function
$f: \set{0,1}^n \rightarrow \set{0,1}$.

\begin{restatable}{theorem}{MainDegree}
  \label{thm:main-degree}
  For all $\eps > 0$ there is a $d = d(\eps)$ such that the following
  holds.  For $n \in \N$, all $s \ge n^{d}$ and any Boolean function
  $f:\set{0,1}^{n} \rightarrow \set{0,1}$ on $n$ bits, SoS requires
  degree $\Omega_\eps(s^{1-\eps})$ to refute $\circuit_s(f)$.
\end{restatable}

It is worthwhile to point out that the proof of \cref{thm:main-degree}
is not specific to the SoS proof system. In fact we outline a general
reduction that shows that if one has a CSP lower bound of the form of
\cref{thm:sos-xor} that only requires good expansion of the underlying
constraint-variable incidence graph and the proof system is expressive
enough so that one can replace variables by local Boolean functions,
then one obtains strong lower bounds for the $\circuit_s(f)$ formula.

The lower bound of $\Omega_\eps(s^{1-\eps})$ on the degree is
essentially tight: if $f$ does not have a circuit of size $s$ then
there exists an SoS refutation of this statement in degree $O(s)$.

\begin{proposition}
  Let $s \in N$ and $f: \set{0,1}^n \rightarrow \set{0,1}$ be a
  Boolean function on $n$ bits that requires circuits of size larger
  than $s$ to be computed. Then there is a degree $O(s)$ SoS
  refutation of $\circuit_s(f)$.
\end{proposition}

We also prove a result about the minimum size (number of monomials)
required for SoS to refute $\circuit_s(f)$. This result holds for all
functions that ``almost'' have a circuit of size $s$, in the sense that
they have an errorless heuristic circuit (see the survey \cite{BT06})
of size $s/2$ and extremely small error probability
with respect to the uniform distribution.
% 
Formally, we let $\calF_n(s,t)$ denote the class of Boolean functions
that consists of all functions $f: \set{0,1}^n \rightarrow \set{0,1}$
for which there is a Boolean circuit
$C_f: \set{0,1}^n \rightarrow \set{0, 1, \bot}$ of size at most $s$
such that
\begin{enumerate}
\item if $C_f(\alpha) \neq \bot$, then $C_f(\alpha) = f(\alpha)$, and
\item $C_f(\alpha) = \bot$ on at most $t$ inputs.
\end{enumerate}
In other words the circuit $C_f$ computes $f$ correctly on all except
$t$ inputs. Note that technically the output of the circuit $C_f$ is
two bits with the first one indicating whether the output is $\bot$ or
the value of the second bit.  We believe that above presentation is
more intuitive and hope that the slight abuse of notation causes no
confusion. With the class of functions $\calF_n(s,t)$ at hand we can
state our main SoS size lower bound.


\begin{restatable}{theorem}{MainSize}
  \label{thm:main-size}
  For all $\eps > 0$ there is a $d = d(\eps)$ such that the following
  holds. Let $n \in N$ and $s \in \N$ such that $s \ge n^d$. If
  $t \ge s$ and $f \in \calF_n(s/2, t)$, then it holds that SoS
  requires size $\exp\big(\Omega_\eps(s^{2-\eps}/t)\big)$ to refute
  $\circuit_s(f)$.
\end{restatable}

This yields non-trivial size lower bounds for $t$ as large as
$s^{2-\eps}/\omega(1)$.  Furthermore, note that once $t \gg s \log s$
there are functions that require such large circuits.  For example
setting $s = 2^{n^{0.99}}$ and $t = s^{1.5}$, the theorem shows that
there are functions $f$ that do not have circuits of size $s$, but SoS
requires size $2^{2^{\Omega(n^{0.99})}}$ to prove this.

It is natural to wonder whether SoS fares better in the monotone
setting. In other words, whether SoS can refute the claim that a
complex monotone function has a small monotone circuit. The following
two theorems show that this is not the case for the set
$\calM_n(\ell)$ of monotone $\ell$-slice functions.  Recall that
$\calM_n(\ell)$ consist of all Boolean functions $f$ on $n$ bits such
that $f(\alpha) = 0$ for all $\alpha$ with Hamming weight less than
$\ell$, and $f(\alpha) = 1$ for all $\alpha$ with Hamming weight
greater than $\ell$ (note that any such $f$ is monotone).

We define a variant $\circuit_s^{\mon}(f)$ of the $\circuit_s(f)$
formula, which instead encodes the claim that $f$ has a monotone
circuit of size $s$, and prove the following theorem.

\begin{restatable}{theorem}{MonotoneDegree}
  \label{thm:main-monotone-degree}
  For all $\eps > 0$ there is a $d = d(\eps)$ such that the following
  holds.  For all $n, \ell \in \N$, all $s \ge n^{d}$ and any monotone
  slice function $f \in \calM_n(\ell)$ SoS requires degree
  $\Omega_\eps(s^{1-\eps})$ to refute
  $\circuit_s^{\mon}(f)$.
\end{restatable}

As in the non-monotone case, we can also obtain size lower bounds for
the monotone-MCSP. Akin to the general size lower bound we consider
monotone Boolean slice functions that have good monotone errorless
heuristic circuits. Let $\calM_n(\ell, s, t) \subseteq \calM_n(\ell)$
be the class of monotone Boolean $\ell$-slice functions
$f: \set{0,1}^n \rightarrow \set{0,1}$ for which there is a
(not necessarily monotone) Boolean circuit
$C_f^\mon: \set{0,1}^n \rightarrow \set{0,1, \bot}$ of size $s$ such
that
\begin{enumerate}
\item for all $\ell$-slice inputs $\alpha \in \binom{[n]}{\ell}$ it
  holds that if $C^\mon_f(\alpha) \neq \bot$, then
  $C^\mon_f(\alpha) = f(\alpha)$, and
\item $C^\mon_f(\alpha) = \bot$ on at most $t$ inputs
  $\alpha \in \binom{[n]}{\ell}$.
\end{enumerate}

\begin{restatable}{theorem}{MonotoneSize}
  \label{thm:main-monotone-size}
  For all $\eps > 0$ there is a $d = d(\eps)$ such that the following
  holds.  For $n, \ell \in \N$, all $s \ge n^{d}$ and $t \ge s$ and
  monotone function $f \in \calM_n(\ell, s/10, t)$ SoS requires size
  $\exp\big(\Omega_\eps(s^{2-\eps}/t)\big)$ to refute
  $\circuit_s^{\mon}(f)$.
\end{restatable}

\subsection{Overview of Proof Techniques}

\paragraph{Degree Lower Bound:}
The main idea that drives our result is a reduction from an expanding
\emph{xor} constraint satisfaction problem to the $\circuit_s(f)$
formula. The reduction is achieved through a careful restriction of
the $\circuit_s(f)$ formula, such that each input
$\alpha \in \set{0, 1}^n$ to the circuit specifies an \emph{xor}
constraint over some new set of variables $Y$.  These $Y$ variables
are a subset of roughly $\Theta(s^{1-\epsilon})$ out of the
$\Theta(s^2)$ many structure variables of the $\circuit_s(f)$ formula.
All other structure variables apart from the $Y$ variables are fixed
to constant values in this step.
This will then result
in an XOR-CSP instance with $2^n$ constraints over the variables
$Y$. All that SoS has to prove is that there is no satisfying
assignment to this XOR-CSP instance. By ensuring that the
constraint-variable incidence graph is sufficiently expanding, SoS
requires large degree to refute the restricted formula (see
\cref{thm:sos-xor}). At the same time, we need the constraint graph to
be very explicit so that it can be encoded into a small circuit.  For
this we utilize a construction of unbalanced expanders by Guruswami et
al.~\cite{GUV09Unbalanced} (see \cref{thm:expander}).  This reduction
then immediately yields \cref{thm:main-degree}.

\pa{Maybe in the above paragraph we also want to point out that all
  the $s \cdot 2^n$ evaluation variables are still alive, effectively
  giving SoS a bunch of auxiliary variables to use when refuting the
  XOR instance, and that some care needs to be taken to deal with
  this?}

\kr{I think this is too technical at this point? I would rather not
  comment on this.}

This lower bound may also be viewed as implementing the general
program sketched by Razborov \cite{Razborov15PRG} relating pseudorandom
generators in proof complexity to the MCSP problem.  However, we prefer
to describe it as a direct reduction to the MCSP problem.

\paragraph{Size Lower Bound:}
In order to obtain size lower bounds, we would like to apply the
degree-size tradeoff due to Atserias and Hakoniemi
\cite{AH18SosTradeoff} to 
\cref{thm:main-degree}. Unfortunately the formula is
over too many variables to be able to conclude a meaningful size lower
bound: it is defined over roughly $\Omega(2^n \cdot s)$
variables.

Instead of applying \cref{thm:main-degree}, we restrict our attention
to functions with all except the at most $t$ $\bot$-outputs computed
by the corresponding errorrless heuristic circuit. If we choose $t$
small enough, then we are able to heavily restrict $\circuit_s(f)$ and
significantly reduce the number of variables to the point where the
Atserias-Hakoniemi degree-size tradeoff is applicable.

\paragraph{Monotone Circuits:}

We prove these theorems by adapting the proofs for the non-monotone setting.
The idea is to work over the $\ell$th slice and disregard all other
inputs. The key feature that makes this work is the fact that the
monotone circuit complexity of a slice function is essentially the same as
the (ordinary) circuit complexity (see \cref{lem:slice}).  This lets us
convert all subcircuits used in the reduction to small monotone circuits (if
we only work on the slice).
% We further need to take care of negations
% in the part of the circuit that computes the \emph{xor}. We do this by
% pushing the negations down until they either hit a $Z$ variable or one
% of the prepared monotone slice circuits.  We create a set
% $\overline Z$ gates, which we can think of as the negation of the
% gates in $Z$ and also make negated versions (on the slice) of the
% monotone circuits that are required for the reduction. This allows us
% to get rid of the last negations in the circuit by appropriately
% connecting the different subcircuits.

The size lower bound goes along the same lines as the proof of
\cref{thm:main-size}. 

\subsection{Organization}

In \cref{sec:preliminaries}, we provide the necessary background
material. In \cref{sec:circ-rest} we set up the general framework for our
lower bounds with some preliminary definitions and lemmas.
Then in \cref{sec:demorgan-LB} we prove the main degree \cref{thm:main-degree} and
size \cref{thm:main-size} lower bounds. We
prove the monotone lower bounds \cref{thm:main-monotone-degree} and
\cref{thm:main-monotone-size} in \cref{sec:monotone-degree-LB}.  In
\cref{sec:upper-bound} we explain how SoS of degree $O(s)$ can refute
$\circuit_s(f)$ (provided $f$ does not have a circuit of size $s$).
Finally in \cref{sec:conclusion} we give some concluding remarks.
